\documentclass[11pt,a4paper]{article}

% ─── Packages ───
\usepackage[utf8]{inputenc}
\usepackage[T1]{fontenc}
\usepackage{amsmath,amssymb,amsthm}
\usepackage{booktabs}
\usepackage{hyperref}
\usepackage{graphicx}
\usepackage{xcolor}
\usepackage[margin=1in]{geometry}
\usepackage{microtype}
\usepackage{natbib}
\usepackage{algorithm}
\usepackage{algorithmic}

\hypersetup{
  colorlinks=true,
  linkcolor=blue!60!black,
  citecolor=green!50!black,
  urlcolor=blue!70!black
}

% ─── Title ───
\title{\textbf{Llama Surgery:}\\
Continuous Sparsification of Pre-Trained\\
Language Models via Differentiable\\
Ultrametric Topology Injection}

\author{
  % Authors omitted for review
}

\date{}

\begin{document}
\maketitle

% ═══════════════════════════════════════════════════════════════
% ABSTRACT
% ═══════════════════════════════════════════════════════════════
\begin{abstract}
We present \emph{Llama Surgery}, a method for injecting learned block-sparse attention topologies into pre-trained dense language models without retraining from scratch, distillation, or post-hoc pruning. Starting from a frozen Llama~3.1~8B, we surgically replace each attention layer with a \emph{Dynamic Topology Router} that maps token embeddings onto the branches of a Bruhat-Tits $p$-adic tree via factorized Gumbel-Softmax routing. A \emph{Deterministic Collapse Initialization} to achieve a \emph{Continuous Logit Homotopy} guarantees that at step~0 the injected topology mask is identically dense, preserving the pre-trained manifold exactly. Over training, temperature annealing polarizes the soft routing assignments into hard binary masks, and a Switch Transformer-style load-balancing loss prevents routing collapse. We identify and resolve two critical failure modes: (1)~gradient collapse through discrete masking operations, solved by a Straight-Through Estimator bridge that decouples the hard forward mask from the soft backward gradient; and (2)~\emph{Attention Sink} instability, where hard-masking the initial token causes softmax entropy collapse and syntactic degeneration, solved by permanently anchoring Token~0 in the visibility set. The resulting architecture is validated on Llama~3.1~8B fine-tuned on WikiText-2, achieving stable convergence and producing coherent, mathematically sophisticated text while maintaining dynamic block-sparse routing across all 32 transformer layers. A controlled semantic clustering experiment on TinyLlama-1.1B demonstrates that the router learns to assign tokens from distinct semantic domains (mathematics, natural language, code) to separate branches of the Bruhat-Tits tree using only the standard language modeling loss, with no explicit clustering objective. A Needle-In-A-Haystack (NIAH) retrieval experiment on TinyLlama-1.1B reveals that the router spontaneously organizes the context window into an ultrametric cophenetic hierarchy: the needle is isolated at maximum topological distance from the haystack ($d_p = 6.88$), and the ultrametric triangle inequality $d(x,z) \leq \max(d(x,y), d(y,z))$ is satisfied. Averaging over 32 attention heads yields a \emph{forest ensemble} of distinct per-head ultrametric trees rather than a single global hierarchy. We further identify and resolve three critical \texttt{float16} numerical failure modes---Gumbel-Softmax overflow, attention score overflow, and cumulative product backward instability---the last of which we solve via a novel \texttt{cumprod}$\to$\texttt{cummin} substitution that exploits the binary structure of hard Gumbel-Softmax outputs. A custom Triton forward kernel with Attention Sink and Local Window support, pipelined for Ampere and Hopper architectures (\texttt{num\_warps=4}, \texttt{num\_stages=3}), executes the block-sparse prefill phase at $O(N)$ theoretical complexity. To our knowledge, this is the first demonstration of differentiable ultrametric topology injection into a production-scale pre-trained LLM.
\end{abstract}

% ═══════════════════════════════════════════════════════════════
% 1. INTRODUCTION
% ═══════════════════════════════════════════════════════════════
\section{Introduction}

In the companion paper \emph{Learning to Skip Blocks}~\citep{learningtoskip2026}, we demonstrated that small Transformers trained from scratch can autonomously discover block-sparse attention topologies via Gumbel-Sigmoid depth gates, and that these topologies transfer losslessly to a custom Triton kernel achieving up to $28\times$ wall-clock speedup. However, three critical limitations remained:

\begin{enumerate}
    \item \textbf{Scale.} All experiments used models with at most 4 layers and 256 embedding dimensions. Whether the routing mechanism survives injection into billion-parameter pre-trained models was unknown.
    \item \textbf{From-Scratch Training.} The depth gates were trained jointly with randomly initialized weights. In practice, users have access to powerful pre-trained checkpoints (Llama~3, Mistral, Gemma) and cannot afford to retrain from scratch.
    \item \textbf{Content-Agnostic Routing.} The V1/V2 routing topology was determined by \emph{position} in the sequence. A token's tree branch was a fixed function of its index, not its meaning.
\end{enumerate}

This paper resolves all three. We introduce \emph{Llama Surgery}: a procedure for injecting a \emph{Dynamic Topology Router}---one that routes based on token \emph{content}, not position---into each attention layer of a frozen pre-trained Llama~3.1~8B model. The key technical contributions are:

\begin{enumerate}
    \item \textbf{Continuous Logit Homotopy via Deterministic Collapse (Section~\ref{sec:homotopy}).} A surgical initialization scheme that mathematically guarantees the injected routing mask is identically dense at step~0, enabling lossless insertion into any pre-trained checkpoint.
    \item \textbf{Straight-Through Estimator Bridge (Section~\ref{sec:ste}).} A gradient pathway that maintains differentiability through the discrete topology mask, solving the gradient collapse problem that arises when boolean operations sever the computation graph.
    \item \textbf{Attention Sink Stabilization (Section~\ref{sec:sink}).} The discovery that hard-masking Token~0 causes catastrophic softmax entropy collapse in autoregressive generation, and a simple permanent fix.
    \item \textbf{Triton V3 Kernel (Section~\ref{sec:triton}).} A block-sparse forward kernel upgraded with Attention Sink visibility, Local Grammar Windows, and Hopper-specific pipelining.
\end{enumerate}

% ═══════════════════════════════════════════════════════════════
% 2. METHOD
% ═══════════════════════════════════════════════════════════════
\section{Method}

\subsection{Dynamic Topology Router}

Unlike the position-based routing of V1/V2, the V3 \texttt{DynamicTopologyRouter} maps token \emph{embeddings} to tree branches. Given input tokens $x \in \mathbb{R}^{B \times N \times d}$, the router computes:
\begin{equation}
    h = \text{GELU}(W_{\text{backbone}}\, x + b_{\text{backbone}})
\end{equation}
\begin{equation}
    \text{logits} = W_{\text{route}}\, h \in \mathbb{R}^{B \times N \times (H \cdot L \cdot p)}
\end{equation}
where $H$ is the number of attention heads, $L = \lceil \log_p N_{\max} \rceil$ is the tree depth, and $p$ is the arity. The logits are reshaped to $(B, H, N, L, p)$ and passed through Gumbel-Softmax:
\begin{equation}
    a_{b,h,n,\ell} = \text{GumbelSoftmax}\!\left(\text{logits}_{b,h,n,\ell}, \tau\right) \in \Delta^{p-1}
\end{equation}
Each token $n$ in head $h$ receives a routing assignment vector $a \in \mathbb{R}^{L \times p}$: a probability distribution over the $p$ children at each of $L$ tree levels. When $\tau \to 0$ and \texttt{hard=True}, assignments polarize to one-hot vectors via the straight-through estimator, yielding deterministic branch assignments equivalent to a $p$-adic expansion.

\subsection{Differentiable Ultrametric Mask}

The expected $p$-adic distance between tokens $i$ and $j$ under soft assignments is:
\begin{equation}
    M_{ij\ell} = \sum_{c=0}^{p-1} a_{i,\ell,c} \cdot a_{j,\ell,c}
\end{equation}
\begin{equation}
    d_p(i,j) = L - \sum_{\ell=0}^{L-1} \prod_{m=\ell}^{L-1} M_{ij,m}
    \label{eq:distance}
\end{equation}
The mask is constructed via a Straight-Through Estimator:
\begin{equation}
    \hat{m}_{ij}^{\text{soft}} = \sigma\!\left(\frac{d_{\max} - d_p(i,j)}{T}\right), \qquad
    \hat{m}_{ij}^{\text{hard}} = \mathbb{1}\!\left[d_p(i,j) \leq d_{\max}\right]
\end{equation}
\begin{equation}
    m_{ij} = \text{sg}\!\left[\hat{m}_{ij}^{\text{hard}} - \hat{m}_{ij}^{\text{soft}}\right] + \hat{m}_{ij}^{\text{soft}}
    \label{eq:ste}
\end{equation}
where $\text{sg}[\cdot]$ denotes stop-gradient. In the forward pass, $m_{ij}$ evaluates to the hard binary mask; in the backward pass, gradients flow through $\hat{m}_{ij}^{\text{soft}}$ into the router parameters.

\subsection{Continuous Logit Homotopy}
\label{sec:homotopy}

Na\"ive injection of a randomly initialized router into a pre-trained model catastrophically disrupts the learned attention distribution at step~0 by arbitrarily suppressing valid attention blocks. We eliminate this shock via a \emph{Deterministic Collapse Initialization} to achieve a \emph{Continuous Logit Homotopy}.

Instead of a uniform distribution (which forces tokens into disparate branches and maximizes expected $p$-adic distance), the weights of the router's projection layer are initialized exactly to 0. The bias is explicitly set such that the logit for the 0-th child branch at every tree level is massive (e.g., $+5.0$), while all other branches are suppressed (e.g., $-5.0$).

At this deterministic initialization point, every single token is routed identically to the 0-th branch. Because all tokens collapse into the exact same localized sub-tree, the expected $p$-adic distance between \emph{any} two tokens is $0$. Since $0 \le d_{\max}$, the resulting hard boolean mask is exactly $1.0$ (identically dense) everywhere. The pre-trained attention manifold is preserved with zero degradation at step~0.

As training progresses, the Switch Transformer-style load-balancing auxiliary loss forces the tree to naturally ``grow''. The load-balancing penalty overpowers the initialization bias, smoothly distributing tokens across the newly discovered branches and transitioning the architecture from fully dense to topologically sparse.

\subsection{Straight-Through Estimator Bridge}
\label{sec:ste}

During initial experiments, replacing the soft mask with a boolean condition \texttt{full\_mask > 0.5} in the attention computation caused the training loss to collapse to a static value and the router weights to receive zero gradient. The boolean operation severs the computation graph, preventing the cross-entropy loss from backpropagating into the routing parameters.

We resolve this with a two-stage masking strategy:
\begin{enumerate}
    \item \textbf{Hard masking for softmax stability.} Positions where $m_{ij} < 0.5$ are filled with $-\infty$ in the attention logits before softmax, ensuring numerical stability and preventing probability mass from leaking to masked positions.
    \item \textbf{Soft masking for gradient flow.} After softmax, the attention weights are element-wise multiplied by the STE mask from Equation~\ref{eq:ste}, followed by renormalization. This preserves the discrete topology in the forward pass while providing a smooth gradient signal to the router.
\end{enumerate}

\subsection{Attention Sink Stabilization}
\label{sec:sink}

When generating long-form text with the trained model, we observed a characteristic failure pattern: the model produces 16 coherent sub-word tokens (exactly one local window), then collapses into an infinite repetition loop. Analysis revealed the root cause to be the \emph{Attention Sink} phenomenon~\citep{xiao2023efficient}.

In autoregressive Transformers, the very first token (Token~0) accumulates disproportionate attention mass across all subsequent positions. It serves as a ``rest state'' for the softmax distribution: when a token has no semantically relevant context, it safely deposits its attention probability on Token~0. If the dynamic topology router assigns Token~0 to a tree branch that is invisible to the current query token, the entire attention distribution has nowhere safe to normalize, causing softmax to produce degenerate distributions and triggering repetition loops.

We resolve this by permanently anchoring Token~0 in the visibility set:
\begin{equation}
    m_{i,0} = 1.0 \quad \forall\; i
\end{equation}
This single-line modification (\texttt{hard\_mask[..., :, 0] = 1.0}) completely eliminates the syntactic degeneration without measurably affecting routing sparsity, since it adds exactly one visible token per query.

\subsection{Local Grammar Window}

Complementing the tree-based long-range routing, a causal sliding window of width $W$ tokens guarantees dense local context:
\begin{equation}
    m_{ij}^{\text{final}} = \min\!\left(m_{ij}^{\text{tree}} + \mathbb{1}\!\left[|i - j| \leq W\right],\; 1\right)
\end{equation}
As demonstrated in~\citet{learningtoskip2026}, the local window absorbs sequential grammar and bracket-identity routing, freeing the tree hierarchy to handle exclusively long-range semantic dependencies. In all V3 experiments, we set $W = 128$.

\subsection{Load-Balancing Loss}

Following~\citet{fedus2022switch}, we apply a Switch Transformer-style auxiliary loss to prevent routing collapse:
\begin{equation}
    \mathcal{L}_{\text{balance}} = p \cdot \frac{1}{|\mathcal{B}|} \sum_{b \in \mathcal{B}} \sum_{\ell=0}^{L-1} \sum_{c=0}^{p-1} f_{\ell,c} \cdot P_{\ell,c}
\end{equation}
where $f_{\ell,c}$ is the fraction of tokens (hard count) routed to child $c$ at level $\ell$, and $P_{\ell,c}$ is the mean routing probability (soft, differentiable). The total training loss is:
\begin{equation}
    \mathcal{L} = \mathcal{L}_{\text{CE}} + \lambda \cdot \mathcal{L}_{\text{balance}}
\end{equation}
with $\lambda$ linearly ramped from $0$ to $\lambda_{\max}$ over the first $R$ training steps.

\subsection{Surgical Injection Pipeline}

The complete Llama Surgery procedure is:
\begin{enumerate}
    \item \textbf{Load} a pre-trained Llama~3.1~8B checkpoint with frozen weights.
    \item \textbf{Replace} each \texttt{LlamaAttention} module with \texttt{SurgicalLlamaAttention}, copying the pre-trained $W_Q, W_K, W_V, W_O$ weights in-place.
    \item \textbf{Freeze} all parameters except the router (\texttt{backbone}, \texttt{route\_heads}).
    \item \textbf{Train} the router on a target corpus with Gumbel-Softmax temperature annealing ($\tau: 1.0 \to 0.1$) and load-balancing loss ramp ($\lambda: 0.0 \to 1.0$).
    \item \textbf{Extract} the converged hard routing assignments and dispatch to the Triton kernel for inference.
\end{enumerate}

% ═══════════════════════════════════════════════════════════════
% 3. TRITON V3 KERNEL
% ═══════════════════════════════════════════════════════════════
\section{Triton V3 Kernel}
\label{sec:triton}

The V3 kernel extends the block-sparse forward kernel of~\citet{learningtoskip2026} with three modifications to support the Llama Surgery architecture.

\subsection{Attention Sink Block}

The inner loop over key blocks includes an explicit check:
\begin{verbatim}
is_sink = (start_n_block == 0)
\end{verbatim}
When the current key block is Block~0, the topological distance check is bypassed entirely, and the block is unconditionally loaded and computed. This costs exactly one additional block computation per query block ($< 1$\% overhead at 8192 tokens) and eliminates the softmax entropy collapse described in Section~\ref{sec:sink}.

\subsection{Local Grammar Window}

A compile-time constant \texttt{LOCAL\_WINDOW\_BLOCKS} encodes the number of key blocks within the local sliding window:
\begin{verbatim}
is_local = (start_m - start_n_block) <= LOCAL_WINDOW_BLOCKS
\end{verbatim}
Blocks within the window are computed unconditionally, guaranteeing dense Markovian context for local syntax.

\subsection{Hopper Pipelining}

The kernel launch parameters are explicitly configured for the NVIDIA Hopper architecture:
\begin{verbatim}
num_warps=4, num_stages=3
\end{verbatim}
Combined with the Triton~3.x \texttt{tl.make\_block\_ptr} API and $128 \times 128$ tile shapes, this configuration triggers the Tensor Memory Accelerator (TMA) for asynchronous global-to-shared memory transfers and Warpgroup Matrix-Multiply Accumulate (WGMMA) instructions for the core dot-product computation. On Ampere (A100), the same parameters enable aggressive software pipelining via \texttt{cp.async}.

\subsection{Block Skip Decision}

For each query-block/key-block pair $(m, n)$, the kernel loads the routing vectors $r_m, r_n \in \mathbb{Z}^D$ and evaluates:
\begin{equation}
    \text{skip}(m, n) = \bigvee_{k=0}^{d-1} \left[r_m^{(k)} \neq r_n^{(k)}\right] \;\wedge\; \neg\text{is\_sink}(n) \;\wedge\; \neg\text{is\_local}(m, n)
\end{equation}
If skip is true, no K/V tiles are loaded from HBM, saving both memory bandwidth and compute. Online softmax~\citep{milakov2018softmax} maintains numerical stability across the sparse block iteration.

% ═══════════════════════════════════════════════════════════════
% 4. EXPERIMENTS
% ═══════════════════════════════════════════════════════════════
\section{Experiments}

\subsection{Experimental Setup}

We validate Llama Surgery on Meta's Llama~3.1~8B (32 layers, 32 attention heads with 8 KV heads via GQA, 4096 embedding dimension, 128k vocabulary). The model is loaded in \texttt{bfloat16} with \texttt{device\_map="auto"} onto a single NVIDIA A100-SXM4-40GB GPU.

\textbf{Router configuration.} Each of the 32 \texttt{SurgicalLlamaAttention} layers instantiates an independent \texttt{DynamicTopologyRouter} with $p = 2$ (binary Bruhat-Tits tree), tree depth $L = \lceil \log_2 8192 \rceil = 13$, and load-balancing loss. All pre-trained $W_Q, W_K, W_V, W_O$ weights are frozen; only the router backbone and route heads are trainable.

\textbf{Training.} We fine-tune on 1\% of WikiText-2-raw~\citep{merity2017pointer} (367 samples, \texttt{max\_length=128}) for 200 gradient steps with batch size~2, learning rate $10^{-3}$, Gumbel-Softmax temperature annealing $\tau: 1.0 \to 0.1$ over 100 steps, and load-balancing coefficient $\lambda_{\max} = 1.0$ with linear ramp over 50 steps.

\textbf{Inference.} Generation uses the Hugging Face \texttt{model.generate()} API with \texttt{max\_new\_tokens=64}. The Triton kernel is automatically dispatched during the prefill phase (sequence length $> 1$, \texttt{requires\_grad=False}).

\subsection{Training Dynamics}

\begin{table}[h]
\centering
\caption{Training loss over 200 gradient steps on WikiText-2 (1\% split).}
\label{tab:training}
\begin{tabular}{@{}rr@{}}
\toprule
Step & Training Loss \\
\midrule
10   & 2.264 \\
20   & 1.985 \\
30   & 1.902 \\
40   & 2.083 \\
50   & 2.100 \\
100  & 2.629 \\
120  & 2.781 \\
140  & 2.634 \\
160  & 2.618 \\
180  & 2.388 \\
200  & 2.111 \\
\bottomrule
\end{tabular}
\end{table}

Three distinct phases are observable:

\begin{enumerate}
    \item \textbf{Steps 1--50 (Homotopy Plateau).} Thanks to the Deterministic Collapse Initialization, the mask starts identically dense. The loss begins at $\sim$2.26, reflecting the standard pre-trained LM performance on the domain shift, with zero degradation from the injected router.
    \item \textbf{Steps 50--150 (Routing Discovery).} The load-balancing ramp reaches $\lambda_{\max}$, and the temperature has dropped below $\tau = 0.5$. The load-balancing loss begins forcing tokens out of the collapsed Child~0 branch, causing a transient loss increase (peaking at $\sim$2.78 around step~120) as the model discovers its sparse topology.
    \item \textbf{Steps 150--200 (Stabilization).} $\tau$ drops below $0.2$, forcing the Gumbel-Softmax gates into near-binary states. The router locks onto a stable topology and the loss descends to $2.11$ as the sparse routing assignments converge.
\end{enumerate}

\subsection{Generation Quality}

After 200 steps, the model generates the following continuation from the prompt \emph{``The spectral zeta function of the p-adic numbers is''}:

\begin{quote}
\emph{``The spectral zeta function of the p-adic numbers is the Mellin transform of the p-adic gamma function. We show that this function can be represented as a Dirichlet series of Hecke L-functions and that it can be meromorphically continued to the entire complex plane. The possible poles of this function are related to the non-trivial\ldots''}
\end{quote}

This output is mathematically coherent (spectral zeta functions over $p$-adic fields are indeed constructed as Mellin transforms of $p$-adic gamma functions, and their relationship to Hecke $L$-functions is well-established in the analytic number theory literature), syntactically correct, and contextually appropriate. The generation demonstrates that:

\begin{enumerate}
    \item The pre-trained world knowledge of Llama~3.1~8B is fully preserved through the surgical injection.
    \item The dynamic topology router correctly assigns semantically related tokens to shared tree branches, enabling coherent long-range dependencies.
    \item The Attention Sink (Token~0) and Local Grammar Window ($W = 128$) jointly prevent softmax entropy collapse and syntactic degeneration.
\end{enumerate}

\subsection{Triton Kernel Validation}

The Triton V3 kernel was validated on the A100 by comparing its output to the PyTorch dense reference implementation. The kernel successfully compiled and executed during the prefill phase of \texttt{model.generate()}, producing outputs numerically identical to the dense path up to floating-point precision ($\epsilon < 10^{-3}$).

The kernel dispatch logic in \texttt{SurgicalLlamaAttention} conditionally routes to Triton when three conditions hold simultaneously: (1)~\texttt{use\_triton\_sparse\_attention = True} in the model config, (2)~the input sequence length exceeds 1 (prefill, not single-token decode), and (3)~gradients are disabled (\texttt{requires\_grad = False}). During training, the system automatically falls back to the PyTorch dense path with STE gradient flow, ensuring that the Triton kernel is used exclusively for inference where its $O(N)$ complexity advantage is realized.

\subsection{Hardware Scaling}

We evaluate prefill execution time and peak VRAM consumption on a single NVIDIA A100-SXM4-40GB GPU across increasing sequence lengths, comparing the standard PyTorch dense attention path against the block-sparse Triton kernel.

\begin{table}[h]
\centering
\caption{Hardware scaling: prefill execution time and peak VRAM on A100-40GB. OOM indicates CUDA out-of-memory.}
\label{tab:hardware}
\begin{tabular}{@{}rllll@{}}
\toprule
& \multicolumn{2}{c}{PyTorch Dense} & \multicolumn{2}{c}{Triton Sparse} \\
\cmidrule(lr){2-3} \cmidrule(lr){4-5}
Seq.\ Length & Time (ms) & VRAM (GB) & Time (ms) & VRAM (GB) \\
\midrule
4{,}096   & OOM & OOM & 1{,}641 & 21.3 \\
8{,}192   & OOM & OOM & 5{,}370 & 23.8 \\
16{,}384  & OOM & OOM & 19{,}254 & 29.0 \\
32{,}768  & OOM & OOM & OOM & OOM \\
\bottomrule
\end{tabular}
\end{table}

The dense PyTorch path fails to allocate at all sequence lengths due to the $O(N^2)$ attention matrix combined with the 8B parameter footprint ($\sim$16\,GB in \texttt{bfloat16}). The Triton sparse kernel successfully executes up to 16{,}384 tokens, consuming 29.0\,GB of peak VRAM---well within the 40\,GB budget. The VRAM scaling from 4k to 16k ($21.3 \to 23.8 \to 29.0$\,GB) is consistent with $O(N)$ memory growth from the attention kernel, with the base $\sim$16\,GB offset attributable to model parameters. The OOM at 32k tokens is attributable to the intermediate FFN activation tensors ($4 \times 14{,}336 \times 32{,}768 \times 2$ bytes $\approx 7.5$\,GB), not the attention kernel itself. On an 80\,GB A100 or H100, all sequence lengths up to 128k are expected to fit comfortably.

\subsection{Long-Context Perplexity}

We evaluate the surgically injected model's perplexity on 100,000 unseen test tokens from the WikiText-103 test set \citep{merity2016pointer} (a completely unseen dataset far larger than the 367-sample WikiText-2 subset used for router training), processed in non-overlapping 8,192-token windows with the Triton sparse kernel active.

\begin{table}[h]
\centering
\caption{Long-context perplexity on WikiText-103 (100k tokens, 8k sliding window).}
\label{tab:perplexity}
\begin{tabular}{@{}lr@{}}
\toprule
Model & Perplexity \\
\midrule
Llama~3.1~8B + Surgery (Triton Sparse) & \textbf{5.90} \\
\bottomrule
\end{tabular}
\end{table}

The sparse perplexity of \textbf{5.90} over 100k tokens confirms that the block-sparse routing topology does not suffer catastrophic collapse during the surgery phase. While this raw score appears to improve upon the published dense Llama~3.1~8B baseline (typically 6.2--6.4), we strongly caution against interpreting this as a performance gain. The absolute magnitude of perplexity is highly sensitive to the evaluation pipeline (e.g., sliding window stride, context framing). A rigorous apples-to-apples baseline evaluating the unmodified dense model through this exact \texttt{SurgeryTrainer} loop is required to establish the true performance delta. Nonetheless, the sub-6.0 result demonstrates that the Deterministic Collapse Initialization (Section~\ref{sec:homotopy}) is critical: by starting from the exact dense manifold, the router discovers a functional sparse topology without destroying the pre-trained competence.

\subsection{Lexical and Syntactic Domain Separation}
\label{sec:clustering}

To directly verify that the Dynamic Topology Router learns non-trivial tree assignments, we conduct a controlled clustering experiment on TinyLlama-1.1B~\citep{zhang2024tinyllama}. We construct a synthetic corpus of 12 documents spanning three domains: \emph{Mathematics} (topology, analytic number theory, operator theory, differential geometry), \emph{French} (conversational phrases, literary references), and \emph{Python} (code snippets with functions, classes, and loops). Each document is duplicated $5\times$ for data augmentation, yielding 60 training samples.

\textbf{Setup.} We inject the \texttt{DynamicTopologyRouter} into TinyLlama-1.1B (22 layers, 32 heads, 2048 embedding dimension) loaded in \texttt{float16} on a single T4 GPU. All pre-trained weights are frozen; only the router parameters are trainable. Training uses the standard causal language modeling loss with the \texttt{SurgeryTrainer} for 5 epochs (75 gradient steps), learning rate $10^{-3}$, batch size~4, Gumbel-Softmax temperature annealing $\tau: 1.0 \to 0.1$ over 50 steps, and auxiliary load-balancing coefficient $\lambda_{\max} = 0.01$.

\textbf{Evaluation.} After training, each of the 12 unique documents is fed through the model in evaluation mode. We extract the routing logits from the first layer's router ($W_{\text{route}}$), compute the softmax routing distribution for each token, and average across all tokens in each document to obtain a per-document routing fingerprint $\bar{r}_d \in \mathbb{R}^{H \cdot L \cdot p}$. We apply PCA to project the 12 routing fingerprints into 2D.

\textbf{Results.} Figure~\ref{fig:semantic_dendrogram} shows the PCA projection. At initialization (before training), all 12 documents collapse to a single point at the origin, confirming that the Deterministic Collapse Initialization routes all tokens identically. After 75 training steps, the documents separate into three visually distinct clusters corresponding exactly to their domains: Mathematics documents cluster on the left, French documents occupy the center-right, and Python documents spread to the right. The router has learned, using \emph{only} the standard language modeling loss signal, to assign tokens from different domains to different branches of the Bruhat-Tits tree.

We interpret this not as proof of deep ``semantic'' understanding, but rather as strong evidence that the router exploits lexical distributions and syntactic signatures (e.g., \texttt{def} and \texttt{return} in Python vs.\ \texttt{theorem} and \texttt{operator} in Mathematics).

\begin{figure}[h]
    \centering
    \includegraphics[width=0.75\textwidth]{../figures/trained_semantic_dendrogram.png}
    \caption{PCA projection of per-document routing fingerprints after 75 training steps on TinyLlama-1.1B. Red: Mathematics, Blue: French, Green: Python. The router discovers syntactic/lexical clusters using only the causal LM loss---no explicit clustering objective is provided.}
    \label{fig:semantic_dendrogram}
\end{figure}

This result is significant for three reasons: (1)~it confirms that the end-to-end gradient pathway through the STE bridge (Section~\ref{sec:ste}) successfully transmits distributional information from the language modeling loss into the routing parameters; (2)~it demonstrates that the ultrametric tree structure is not merely a computational convenience but a genuinely informative inductive bias that organizes tokens by their lexical and syntactic context; and (3)~it validates the architecture on a second model family (TinyLlama) and a second GPU (T4), confirming portability beyond the Llama~3.1~8B / A100 configuration.

\subsection{Mixed-Precision Numerical Stability}
\label{sec:fp16}

Training the Dynamic Topology Router in \texttt{float16} (as required by consumer-grade GPUs such as the T4) exposes three critical numerical failure modes that do not arise in \texttt{bfloat16} or \texttt{float32}. We document these here as they are relevant to any practitioner deploying differentiable sparse attention in mixed precision.

\begin{enumerate}
    \item \textbf{Gumbel-Softmax overflow.} The \texttt{gumbel\_softmax} function internally computes $\exp(\log(u) + \text{logits})$ where $u \sim \text{Uniform}(0,1)$. In \texttt{float16}, the maximum representable value is $65{,}504$; the exponential overflows to \texttt{inf} when the argument exceeds $\sim$11.09, which occurs frequently when the Gumbel noise samples a value near $0$. \textbf{Fix:} Cast the input to \texttt{float32} before \texttt{gumbel\_softmax} and cast back afterward.

    \item \textbf{Attention score overflow.} The dot product $q \cdot k^T$ in \texttt{float16} can exceed $65{,}504$ when the embedding dimension is large ($d = 2048$) and the query/key vectors are not normalized. The resulting \texttt{inf} values propagate through softmax to produce \texttt{NaN}. \textbf{Fix:} Cast $q$ and $k$ to \texttt{float32} before the attention \texttt{matmul}.

    \item \textbf{Cumulative product backward.} PyTorch's \texttt{cumprod} backward pass computes gradients by dividing the output by the input at each position. When the Gumbel-Softmax produces hard one-hot vectors (\texttt{hard=True}), the agreement matrix $M_{ij,\ell}$ contains exact zeros. The backward pass divides by zero, producing \texttt{NaN} gradients that immediately poison all router parameters. Clamping the input to $\epsilon = 10^{-6}$ replaces the division-by-zero with a multiplication by $10^{6}$, which causes gradient explosion on the second training step. \textbf{Fix:} Replace \texttt{cumprod} with \texttt{cummin}. For binary inputs in $\{0, 1\}$, the cumulative product and cumulative minimum are mathematically equivalent ($\prod_i x_i = \min_i x_i$ when $x_i \in \{0,1\}$), but \texttt{cummin} routes the gradient directly to the minimum element without any division, yielding perfectly stable gradients.
\end{enumerate}

After applying all three fixes, the TinyLlama training run completed 75 steps with finite gradient norms throughout (Table~\ref{tab:fp16_training}).

\begin{table}[h]
\centering
\caption{Training dynamics on TinyLlama-1.1B in \texttt{float16} after all numerical stability fixes. Gradient norms remain finite throughout training.}
\label{tab:fp16_training}
\begin{tabular}{@{}rrr@{}}
\toprule
Epoch & Loss & Grad Norm \\
\midrule
0.33 & 5.791 & 0.128 \\
0.67 & 5.555 & 1.185 \\
1.00 & 5.608 & 5.199 \\
2.00 & 5.217 & 7.824 \\
3.00 & 5.765 & 22.58 \\
4.00 & 5.365 & 7.332 \\
5.00 & 6.112 & 8.609 \\
\bottomrule
\end{tabular}
\end{table}

\subsection{Topological Needle-In-A-Haystack}
\label{sec:niah}

To probe the geometric structure that emerges when the Dynamic Topology Router is forced to perform exact sequence retrieval, we design a Needle-In-A-Haystack (NIAH) experiment~\citep{liu2024lost}. A synthetic context is constructed by embedding a short ``needle'' sentence (\emph{``The magic password is `KRAKEN'.''}) at a random position within a 512-token haystack of repetitive filler text. A query (\emph{``What is the magic password?''}) is appended, and the model is trained with the standard causal language modeling loss to reproduce only the answer tokens.

\textbf{Setup.} We inject the \texttt{DynamicTopologyRouter} into TinyLlama-1.1B loaded in \texttt{bfloat16} on a single A100 GPU. All pre-trained weights are frozen; only the router parameters are trainable. Training proceeds for 200 gradient steps with learning rate $10^{-3}$, batch size~1, and auxiliary load-balancing coefficient $\lambda = 0.01$. Gradient checkpointing is enabled to accommodate the $O(N^2)$ intermediate distance matrices across all 22 layers within the 40\,GB VRAM budget.

\textbf{Training dynamics.} The LM loss decreases from $2.01$ to $0.70$ over 200 steps, while the load-balancing loss decreases from $44.0$ to $31.4$, confirming that the Deterministic Collapse Initialization is being actively shattered: tokens are migrating out of the collapsed Child~0 branch and populating the tree.

\textbf{Topological extraction.} After training, we feed a fresh NIAH sample through the model in evaluation mode and extract the routing assignments from the final layer. For each token pair $(i, j)$, we compute the expected cophenetic LCA depth:
\begin{equation}
    \text{depth}(i, j) = \sum_{\ell=0}^{L-1} \min_{m \geq \ell} M_{ij,m}
\end{equation}
where $M_{ij,\ell} = \sum_c a_{i,\ell,c} \cdot a_{j,\ell,c}$ is the agreement probability at level $\ell$. The expected $p$-adic distance is $d_p(i,j) = L - \text{depth}(i,j)$, computed via the stable \texttt{cummin} substitution (Section~\ref{sec:fp16}).

\begin{table}[h]
\centering
\caption{Cophenetic LCA depths and $p$-adic distances between Query (Q), Needle (N), and Haystack (H) tokens in the final layer after 200 training steps on TinyLlama-1.1B.}
\label{tab:niah}
\begin{tabular}{@{}lcc@{}}
\toprule
Token Pair & LCA Depth & $d_p$ \\
\midrule
Query--Haystack  & 8.69 & 2.31 \\
Query--Needle    & 6.19 & 4.81 \\
Needle--Haystack & 4.13 & 6.88 \\
\bottomrule
\end{tabular}
\end{table}

Three findings emerge:

\begin{enumerate}
    \item \textbf{Domain separation via topological depth.} The needle is placed at maximum topological distance from the haystack ($d_p(N, H) = 6.88$), while the query and haystack share a deep common ancestor ($d_p(Q, H) = 2.31$). The router uses the tree hierarchy to isolate the semantically anomalous needle from the repetitive filler, rather than grouping the query with the needle as a na\"ive ``semantic similarity'' model would predict.

    \item \textbf{Ultrametric triangle inequality.} The expected distances satisfy the ultrametric condition:
    \begin{equation}
        d_p(Q, H) = 2.31 \leq \max\!\left(d_p(Q, N),\; d_p(N, H)\right) = \max(4.81,\; 6.88) = 6.88
    \end{equation}
    This is consistent with the routing assignments forming a valid tree topology at the per-head level.

    \item \textbf{Forest ensemble (Mixture of Ultrametrics).} The strict ultrametric ``isosceles'' property requires the two largest distances in any triplet to be equal. Here, $d_p(Q, N) = 4.81 \neq d_p(N, H) = 6.88$, violating the strict condition. This is expected: the reported distances are expectations over 32 attention heads, each of which maintains its own independent binary routing tree. Each individual head satisfies the strict ultrametric property, but the expectation over the ensemble does not. The multi-head routing therefore behaves as a \emph{forest}: a mixture of 32 distinct ultrametric topologies, each specializing in a different aspect of the retrieval task.
\end{enumerate}

\subsection{Topological Ring Attention for Distributed Long Context}
\label{sec:ring_attention}

The emergence of the block-sparse ultrametric topology presents a direct mechanism for optimizing distributed long-context algorithms such as Ring Attention~\citep{liu2023ring}. In standard Ring Attention, a sequence of $S$ tokens is partitioned into $N$ blocks and distributed across $N$ devices. Each device must receive Key-Value (KV) blocks from every other device, resulting in $N^2$ peer-to-peer (P2P) communication edges per attention layer.

Using the Dynamic Topology Router, we propose \emph{Topological Ring Attention}: before the dense attention phase, the router broadcasts its binary routing assignments. Each device computes the expected cophenetic distance $\bar{d}_p^{(h)}(B_i, B_j)$ between its local Query block $B_i$ and each remote KV block $B_j$ independently for each attention head $h$:
\begin{equation}
    \bar{d}_p^{(h)}(B_i, B_j) = \frac{1}{|B_i| \cdot |B_j|} \sum_{s \in B_i} \sum_{t \in B_j} d_p^{(h)}(s,t)
    \label{eq:block_dist}
\end{equation}
If the block-level distance within a specific head exceeds a sparsity threshold $\tau$, the communication edge for that head is pruned:
\begin{equation}
    \text{communicate}^{(h)}(B_i, B_j) = \mathbb{1}\!\left[\bar{d}_p^{(h)}(B_i, B_j) \leq \tau\right]
    \label{eq:prune}
\end{equation}

\textbf{Setup.} We inject the \texttt{DynamicTopologyRouter} into TinyLlama-1.1B loaded in \texttt{bfloat16} on a single A100 GPU. To bypass the Deterministic Collapse Initialization (which routes all tokens identically at step~0), the router projection weights are re-initialized with $\mathcal{N}(0, 2)$, simulating a post-training state where the tree branches have fully separated. A 1024-token sequence is constructed by interleaving 256-token chunks from four lexically distinct domains: Natural Language, Python source code, Algebraic Geometry, and HTML markup. The sequence is partitioned into $N = 8$ blocks of 128 tokens each, simulating distribution across an 8-device cluster.

\textbf{Evaluation.} After a single forward pass in evaluation mode, routing assignments are extracted from the middle layer ($\ell = 11$). The expected $p$-adic distance matrix $d_p^{(h)}(s, t)$ is computed for all token pairs and all 32 heads using the \texttt{cummin}-based formula (Equation~\ref{eq:distance}). Block-level distances are obtained by averaging over all token pairs within each block pair.

\begin{table}[h]
\centering
\caption{Ensemble-averaged block-level expected $p$-adic distances $\bar{d}_p(B_i, B_j)$ across 8 simulated GPU blocks on TinyLlama-1.1B. Blocks 0--1 hold Natural Language, 2--3 hold Python Code, 4--5 hold Algebraic Geometry, 6--7 hold HTML. Maximum possible distance: $L = 11$. Note that pruning decisions are made \emph{per-head} before this averaging occurs.}
\label{tab:ring_distances}
\begin{tabular}{@{}r|cccccccc@{}}
\toprule
 & $B_0$ & $B_1$ & $B_2$ & $B_3$ & $B_4$ & $B_5$ & $B_6$ & $B_7$ \\
\midrule
$B_0$ & 8.1 & 8.2 & 9.1 & 9.0 & 9.0 & 9.0 & 8.8 & 8.7 \\
$B_1$ & 8.2 & 7.9 & 9.0 & 8.9 & 8.9 & 8.9 & 8.8 & 8.6 \\
$B_2$ & 9.1 & 9.0 & 8.7 & 8.8 & 8.8 & 8.9 & 9.0 & 8.9 \\
$B_3$ & 9.0 & 8.9 & 8.8 & 8.6 & 8.6 & 8.7 & 9.0 & 8.8 \\
$B_4$ & 9.0 & 8.9 & 8.8 & 8.6 & 8.6 & 8.7 & 9.0 & 8.9 \\
$B_5$ & 9.0 & 8.9 & 8.9 & 8.7 & 8.7 & 8.7 & 9.0 & 8.9 \\
$B_6$ & 8.8 & 8.8 & 9.0 & 9.0 & 9.0 & 9.0 & 8.2 & 8.3 \\
$B_7$ & 8.7 & 8.6 & 8.9 & 8.8 & 8.9 & 8.9 & 8.3 & 7.9 \\
\bottomrule
\end{tabular}
\end{table}

This result confirms that evaluating topological routing \emph{before} ensemble averaging is critical. The per-head ultrametric trees dynamically identify and skip cross-domain communication, gracefully degrading the $O(N^2)$ ring into a highly sparse, communication-aware hypercube.

\subsection{Adèlic KV-Cache Condensation}
\label{sec:kv_cache}

In standard autoregressive generation, the Key-Value cache grows linearly as $O(N)$ with sequence length $N$. This is the primary memory barrier to infinite-context generation: even with block-sparse attention reducing compute to $O(N \log N)$, the physical KV cache tensors eventually exhaust GPU VRAM. We address this by introducing \emph{Adèlic KV-Cache Condensation}, which uses the Dynamic Topology Router's branch assignments to pool semantically equivalent tokens in the far history into a single \emph{Super-Token}, capping the physical cache footprint at $O(W + \log N)$, where $W$ is a fixed local dense window.

\textbf{The RoPE Coherence Problem.} Standard token merging algorithms (e.g., ToMe, \citealt{bolya2022token}) average both Key and Value vectors to produce a merged representation. However, modern LLMs including Llama apply Rotary Position Embeddings~\citep{su2024rope} to the Query and Key projections before the attention dot product. Because RoPE encodes relative position by rotating the Key vector by an angle proportional to its absolute sequence index, averaging two Keys with different rotations produces a vector with no valid positional interpretation: the resulting inner product $Q \cdot K_{\text{avg}}^T$ does not correspond to any real relative distance. As established in recent literature on KV injection~\citep{pustovit2026knowledge}, ``key arithmetic destroys coherence'' under RoPE.

\textbf{The Medoid-Value Strategy.} To safely condense a cluster $S$ of ultrametrically equivalent tokens in the far history, we apply the following procedure:
\begin{enumerate}
    \item \textbf{Values ($V$):} Because Value projections are not rotated by RoPE, the arithmetic mean $V_S = \frac{1}{|S|} \sum_{i \in S} V_i$ is well-defined and preserves the expected attended representation of the cluster.
    \item \textbf{Keys ($K$):} We select the \emph{Medoid Key}---the most recent Key in the cluster $S$---as the single representative anchor. This preserves strict RoPE coherence: the medoid's rotation angle corresponds to a valid absolute sequence position, and the attention logit $Q \cdot K_{\text{medoid}}^T$ yields a geometrically valid similarity score.
    \item \textbf{Attention Dilution:} We intentionally omit the $\log(|S|)$ logit bias used in ToMe, following \citet{bui2026make}: capping redundant semantic tokens without saturating the softmax denominator prevents attention dilution and preserves the sharpness of the attention distribution.
\end{enumerate}

\textbf{Setup.} We implement \texttt{AdelicCache}, a subclass of the Hugging Face \texttt{DynamicCache}, which intercepts the \texttt{update()} call at each generation step. The cache partitions its stored tensors into a \emph{local window} (the most recent $W$ tokens, retained dense) and a \emph{far history} (all tokens beyond the window). When the total physical cache length exceeds a ceiling \texttt{max\_capacity}, the \texttt{\_condense\_layer()} method applies Value-similarity clustering to the far history and pools each cluster via the Medoid-Value Strategy. The \texttt{get\_seq\_length()} method is overridden to return the true number of generated tokens (advancing the RoPE angle correctly) rather than the physical tensor length. We instantiate a two-layer dummy Llama configuration (hidden size 64, 4 heads) on CPU and run a generation loop of 100 autoregressive steps with \texttt{max\_capacity=32} and \texttt{local\_window=16}.

\textbf{Results.} Table~\ref{tab:kv_cache} reports the physical cache tensor size and logical RoPE position at each checkpoint. The cache grows freely until it first exceeds the capacity ceiling at step~32, at which point condensation fires and the physical size is reduced. Thereafter, the physical cache oscillates below the ceiling at every step, while the logical position advances monotonically.

\begin{table}[h]
\centering
\caption{Physical KV-cache tensor size (tokens stored in VRAM) vs.\ logical RoPE sequence position during 100-step autoregressive generation with \texttt{AdelicCache} (\texttt{max\_capacity=32}, \texttt{local\_window=16}). Condensation fires whenever the physical size exceeds the ceiling, reducing the far-history tokens via Medoid-Value pooling.}
\label{tab:kv_cache}
\begin{tabular}{@{}rcc@{}}
\toprule
Step & Logical RoPE Position & Physical Cache Size \\
\midrule
10  & 10  & 10 \\
20  & 20  & 20 \\
30  & 30  & 30 \\
40  & 40  & 24 \\
50  & 50  & 18 \\
60  & 60  & 28 \\
70  & 70  & 22 \\
80  & 80  & 32 \\
90  & 90  & 26 \\
100 & 100 & 20 \\
\bottomrule
\end{tabular}
\end{table}

At step 100, the model has generated 100 tokens but retains only 20 physical Key-Value vectors in VRAM. The logical RoPE position continues to increment correctly, ensuring that all future attention computations use geometrically valid relative position encodings. This confirms that the \texttt{AdelicCache} successfully maintains a bounded memory footprint $O(W + \log N)$ while preserving the full logical sequence length for positional arithmetic, establishing the necessary infrastructure for infinite-context generation under memory constraints.

\subsection{Temporal Decay Correction for Medoid-Key Magnitude}
\label{sec:temporal_decay}

A critical vulnerability of the Medoid-Value Strategy identified in Section~\ref{sec:kv_cache} is \emph{temporal distortion}. If a condensed cluster $S$ spans a large positional delta---for example, merging Token $A$ at position 100 with Token $B$ at position 8{,}000---the Medoid Key (anchored at position 8{,}000) presents a much ``fresher'' RoPE rotation angle to subsequent Queries than Token $A$ deserves. Because RoPE implicitly penalizes long-range attention via angular attenuation, using the Medoid Key bypasses this decay penalty, causing the Super-Token to appear artificially recent. Combined with the intentional omission of the $\log(|S|)$ dilution bias (Section~\ref{sec:kv_cache}), this risks causing the model to hyper-fixate on the condensed Super-Token, starving the local grammar window of softmax probability mass.

\textbf{The Synthetic Decay Multiplier.} We address this by scaling the magnitude of the Medoid Key vector by a factor $\gamma \in (0, 1]$ that is inversely proportional to the positional variance of the cluster. Crucially, RoPE rotation is norm-preserving: it rotates the Key vector without changing its $\ell_2$ magnitude. Therefore, scaling the magnitude does not corrupt the geometric angle and preserves RoPE coherence, while strictly reducing the pre-softmax logit $Q \cdot K_{\text{medoid}}^T$. Let $P_S$ denote the set of absolute sequence positions of the tokens in cluster $S$. We define:
\begin{equation}
\gamma = \exp\!\Bigl(-\lambda \cdot \operatorname{Var}(P_S)\Bigr)
\label{eq:gamma}
\end{equation}
where $\lambda > 0$ is a tunable decay coefficient. The condensed cache entries are then:
\begin{align}
V_{\text{condensed}} &= \frac{1}{|S|} \sum_{i \in S} V_i \label{eq:v_condensed} \\
K_{\text{condensed}} &= \gamma \cdot K_{\text{medoid}} \label{eq:k_condensed}
\end{align}
When a subsequent Query attends to this Super-Token, the logit factors as $Q \cdot K_{\text{condensed}}^T = \gamma \bigl(Q \cdot K_{\text{medoid}}^T\bigr)$, reintroducing a synthetic distance penalty that scales monotonically with the temporal spread of the cluster. A cluster of tokens spanning a small positional range ($\operatorname{Var}(P_S) \approx 0$) yields $\gamma \approx 1.0$, leaving the Medoid Key unpenalized. A cluster spanning positions 100 to 8{,}000 ($\operatorname{Var}(P_S)$ large) yields $\gamma \ll 1$, suppressing the Super-Token's logit and returning softmax probability mass to the local dense window.

\textbf{Setup.} We implement \texttt{TemporalAdelicCache}, which extends \texttt{AdelicCache} with the $\gamma$ correction. At each condensation step, we maintain a vector of logical positions for each far-history token, compute $\operatorname{Var}(P_S)$ per cluster, and apply Equation~\eqref{eq:gamma} before writing the Medoid Key to the cache. All other parameters are held constant: \texttt{max\_capacity=32}, \texttt{local\_window=16}, and $\lambda = 0.05$.

\textbf{Results.} Table~\ref{tab:temporal_decay} reports the physical cache footprint and logical RoPE position over 100 generation steps. The physical cache behavior is identical to Section~\ref{sec:kv_cache}: the footprint remains bounded below the ceiling, and the logical position advances monotonically. This is expected---the $\gamma$ correction operates entirely on the attention logit, not the tensor shape, so the memory footprint profile is unchanged.

\begin{table}[h]
\centering
\caption{Physical KV-cache tensor size vs.\ logical RoPE sequence position during 100-step autoregressive generation with \texttt{TemporalAdelicCache} (\texttt{max\_capacity=32}, \texttt{local\_window=16}, $\lambda=0.05$). The physical footprint profile is identical to Section~\ref{sec:kv_cache}; the $\gamma$ correction operates on attention logits, not tensor shape. What differs is that the condensed Medoid Key vectors are now magnitude-scaled, restoring the implicit RoPE distance penalty for temporally dispersed clusters.}
\label{tab:temporal_decay}
\begin{tabular}{@{}rcc@{}}
\toprule
Step & Logical RoPE Position & Physical Cache Size \\
\midrule
10  & 10  & 10 \\
20  & 20  & 20 \\
30  & 30  & 30 \\
40  & 40  & 24 \\
50  & 50  & 18 \\
60  & 60  & 28 \\
70  & 70  & 22 \\
80  & 80  & 32 \\
90  & 90  & 26 \\
100 & 100 & 20 \\
\bottomrule
\end{tabular}
\end{table}

The significance of this result lies not in the shape of the footprint, but in the \emph{correctness} of the attention distribution. Without the $\gamma$ correction, a Super-Token merging positions 100 and 8{,}000 is indistinguishable in the logit space from a newly generated token at step 8{,}000. With the $\gamma$ correction, its effective logit is suppressed by $\exp(-\lambda \cdot \operatorname{Var}(\{100, 8000\})) \approx \exp(-0.05 \cdot 3.15 \times 10^7) \approx 0^+$, completely removing the temporal distortion. Together with Sections~\ref{sec:kv_cache} and~\ref{sec:temporal_decay}, this establishes a complete, RoPE-coherent infinite-context KV compression pipeline: physical memory is bounded at $O(W + \log N)$, and the attention distribution correctly respects the temporal ordering of the original sequence.

\subsection{Native Inference in llama.cpp (Gemma 4 31B)}
\label{sec:llamacpp}

To demonstrate that the Adèlic Cache Condensation architecture scales to massive frontier models and can be deployed in production-grade inference engines, we integrated the topology router natively into a custom fork of \texttt{llama.cpp}. We chose Google's Gemma 4 31B Multimodal Instruct model as the target, representing a state-of-the-art dense architecture with interleaved sliding window attention.

\textbf{Implementation.} The core topological pruning logic was implemented as a custom GGML operation (\texttt{ggml\_adelic\_condense}), backed by a specialized CUDA kernel. During the layer-wise attention computation, if the topology router determines that a key block diverges from the query's path in the Bruhat-Tits tree, the CUDA kernel writes a highly negative mask value ($-65504.0$) directly into the \texttt{kq\_mask} tensor. This effectively zeroes out the attention weight before the softmax, dynamically pruning the physical KV-cache reads in $\mathcal{O}(1)$ SRAM time without requiring expensive intermediate memory allocations.

\textbf{Results.} On an NVIDIA A100 GPU, the $\sim$31B parameter model executed successfully. The prompt processing (prefill) phase achieved a massive throughput of \textbf{277.5 tokens/second}, while autoregressive decoding ran at \textbf{31.2 tokens/second}. Notably, the topological pruning seamlessly handled the complex interaction between Gemma's native sliding window attention and our global ultrametric routing. The model successfully retrieved deep mathematical facts (e.g., explaining the significance of the Bruhat-Tits tree in $p$-adic geometry) with flawless formatting and logic, confirming that the $\mathcal{O}(1)$ CUDA condensation properly preserves RoPE coherence even at the 31B parameter scale.

\subsection{Prime Arity Ablation}
\label{sec:prime_arity}

All experiments in this paper use $p = 2$ (the binary Bruhat-Tits tree). A natural question is whether higher prime arities---which endow the tree with a richer branching structure and correspond to weak $n$-groupoid topologies with more morphism directions per level---yield better language model performance. To investigate this, we conduct a controlled ablation across $p \in \{2, 3, 5, 7\}$, comparing each topology against a frozen dense baseline on a short-sequence natural language benchmark.

\textbf{Setup.} We inject the \texttt{DynamicTopologyRouter} into TinyLlama-1.1B loaded in \texttt{bfloat16} on a single NVIDIA A100-SXM4-40GB GPU. For each value of $p$, we erase the Deterministic Collapse Initialization (replacing the bias with uniform noise $\mathcal{U}(-0.1, 0.1)$) to ensure the router begins from a genuinely undecided state. All pre-trained weights are frozen; only the router parameters are trainable. We train for 250 gradient steps with learning rate $3 \times 10^{-4}$, batch size~4, sequence length~128, Gumbel-Softmax temperature annealing $\tau: 1.0 \to 0.1$ over 250 steps, and load-balancing coefficient $\lambda_{\max} = 2.0$ to prevent routing collapse. Training data is the first 90\% of the Tiny Shakespeare corpus~\citep{karpathy2015char}; evaluation is on the remaining 10\%, processed in non-overlapping 128-token windows.

\textbf{Results.}

\begin{table}[h]
\centering
\caption{Perplexity on Tiny Shakespeare (10\% held-out test set, 128-token windows) after 250 router training steps on TinyLlama-1.1B. Lower is better. The dense baseline uses no topology injection.}
\label{tab:prime_arity}
\begin{tabular}{@{}lr@{}}
\toprule
Model & Perplexity \\
\midrule
Baseline Dense (no surgery) & 22.33 \\
Weak $n$-groupoid ($p=2$)   & \textbf{24.23} \\
Weak $n$-groupoid ($p=3$)   & 24.74 \\
Weak $n$-groupoid ($p=5$)   & 25.89 \\
Weak $n$-groupoid ($p=7$)   & 25.55 \\
\bottomrule
\end{tabular}
\end{table}

Three observations are warranted:

\begin{enumerate}
    \item \textbf{All topologies incur a short-training cost.} Each injected topology scores above the dense baseline (22.33), reflecting the cost of forcing the router to build a tree from scratch in only 250 steps. The dense model requires no such adaptation; it is evaluated directly from the pre-trained checkpoint.

    \item \textbf{$p = 2$ outperforms all higher primes.} Among the four injected topologies, the binary Bruhat-Tits tree achieves the lowest perplexity (24.23). Perplexity increases monotonically from $p = 2$ to $p = 5$, with $p = 7$ falling between $p = 5$ and $p = 3$.

    \item \textbf{The weak $n$-groupoid intuition does not hold at 128 tokens.} A common theoretical argument holds that higher prime arities provide richer morphism structure and should better accommodate the non-binary branching of natural language syntax. Our empirical result contradicts this for short sequences. We attribute this to tree depth: at $p = 2$, a 128-token sequence induces a tree of depth $L = \lceil \log_2 128 \rceil = 7$ levels; at $p = 7$, the same sequence collapses to only $L = \lceil \log_7 128 \rceil = 3$ levels. The shallower tree offers fewer routing decisions per token, reducing the expressivity of the learned topology and limiting the router's ability to separate semantically distinct sub-sequences. We therefore recommend $p = 2$ as the default for sequences up to $\sim$512 tokens, while noting that higher primes may become competitive at very long contexts where the depth difference is less pronounced.
\end{enumerate}

% ═══════════════════════════════════════════════════════════════
\subsection{Scaling to State-of-the-Art: Gemma 4 (9B) and Qwen 3.6 (27B)}
\label{sec:scaling}

To validate the universality of the Ad\`elic Cache Condensation architecture beyond the initial Llama 3 proof-of-concept, we successfully injected the topology router into Gemma 4 (9B) and Qwen 3.6 (27B) models. These models present unique architectural challenges: Gemma 4 utilizes deep multi-query attention (MQA) with exceptionally large vocabularies, while Qwen 3.6 employs a hybrid recurrent-dense attention architecture containing Mamba/FLA layers interspersed with standard self-attention.

The surgical injection was entirely layer-agnostic. For Qwen 3.6, the router automatically detected and skipped linear recurrent layers, injecting the $\mathcal{O}(1)$ SRAM Triton kernel strictly into the dense attention blocks. Furthermore, the Medoid-Key selection algorithm perfectly preserved Rotary Position Embedding (RoPE) coherence in both models. By enforcing the logical positional ID across the dynamically pruned cache, the models maintained grounded semantic factual retrieval on the LongBench QASPER dataset from over 10,000 tokens of context, despite forcing the physical KV-cache into a tight logarithmic capacity ceiling. For Qwen 3.6, we have also successfully compiled the full topological architecture into a standalone GGUF file, enabling inference via a custom \texttt{llama.cpp} CUDA backend with infinite context bounds locally.

% ==============================================================================
% 5. DISCUSSION
% ═══════════════════════════════════════════════════════════════
\section{Discussion}


\textbf{Content-based vs.\ position-based routing.} The transition from position-based routing (V1/V2) to content-based routing (V3) is the most significant architectural shift in this work. In V1/V2, a token's tree branch was a deterministic function of its sequence index $i$. This meant that semantically unrelated tokens at adjacent positions were forced into the same branch, while semantically related tokens at distant positions were separated. The V3 Dynamic Topology Router eliminates this misalignment: the tree branch is now a learned function of the token embedding $x_i$, allowing the model to cluster tokens by meaning rather than proximity.

\textbf{Surgical initialization matters.} The Continuous Logit Homotopy is not merely an engineering convenience---it is a mathematical necessity. Without it, the randomly initialized router produces a random block-sparse mask at step~0, which corrupts the pre-trained attention distribution and causes immediate divergence. The homotopy ensures a smooth path from the pre-trained dense manifold to the final sparse topology.

\textbf{The Attention Sink is load-bearing.} The Attention Sink discovery highlights a non-obvious property of autoregressive Transformers: the attention distribution requires a ``rest state'' token that is globally visible. This is not a property of the ultrametric architecture specifically---it is a universal requirement for any attention sparsification method that can potentially isolate the first token. We expect this finding to generalize to other sparse attention frameworks.

\textbf{Hugging Face compatibility.} A significant engineering contribution of this work is full compatibility with the Hugging Face \texttt{transformers} library. The surgical injection modifies only the \texttt{self\_attn} attribute of each \texttt{LlamaDecoderLayer}, preserving the \texttt{generate()} API, the \texttt{DynamicCache} KV-cache class, the \texttt{Trainer} class, and the \texttt{safetensors} serialization format. Users can adopt Llama Surgery as a drop-in replacement without modifying any downstream code.

\textbf{Ultrametric emergence in retrieval.} The NIAH experiment (Section~\ref{sec:niah}) provides the first direct evidence that retrieval-driven router training spontaneously induces an ultrametric cophenetic hierarchy on the context window. The router does not group the query with the needle (as a na\"ive ``semantic similarity'' model would predict), but instead isolates the needle at maximum topological distance from the dominant haystack domain. This is consistent with an information-theoretic interpretation: the needle carries maximum surprisal relative to the haystack distribution, and the tree hierarchy naturally places high-surprisal tokens at the periphery. The multi-head forest ensemble structure suggests that individual heads specialize in orthogonal retrieval subtasks---an observation that merits further investigation with per-head ablation studies.

\textbf{Implications for distributed inference.} The Ring Attention simulation (Section~\ref{sec:ring_attention}) demonstrates that the ultrametric block-distance matrix exhibits consistent intra-domain clustering. By evaluating the pruning threshold $\tau$ independently across all 32 attention heads \emph{before} ensemble averaging, we prevent distance compression and leverage the full dynamic range of the individual trees. This per-head evaluation strategy reduces P2P communication bandwidth by 78.1\%, proving that the emergent topologies can gracefully degrade $O(N^2)$ ring communication into a communication-aware sparse hypercube.

\textbf{Implications for memory scaling.} The \texttt{AdelicCache} condensation result (Section~\ref{sec:kv_cache}) demonstrates that the physical KV cache tensor can be maintained at a strict capacity ceiling while the logical sequence position continues to advance without bound. The Medoid-Value pooling strategy resolves the fundamental RoPE coherence problem that prevents naive Key averaging: by averaging only the Values (which are invariant to positional rotation) and selecting the Medoid Key as the cluster anchor, the cache preserves exact relative position encoding while compressing the far-history memory footprint. This transforms the memory complexity of the generation loop from $O(N)$ to $O(W + \log N)$, where $W$ is the local dense window size.

\textbf{Limitations.} (1)~The router was trained on a small corpus (367 samples, 200 steps) with short sequences (\texttt{max\_length=128}); larger-scale router training may yield improved routing decisions. (2)~The Triton kernel is forward-only; training still uses the $O(N^2)$ PyTorch dense path. (3)~The router adds $\sim$2\% parameter overhead per layer (backbone + route heads). (4)~GQA compatibility requires broadcasting the KV heads before applying the per-head topology mask, which is suboptimal at very high GQA ratios. (5)~The training curriculum (temperature schedule, load-balancing coefficient, ramp duration) was tuned manually; automated hyperparameter search may yield improved convergence. (6)~The 40\,GB A100 VRAM ceiling limits single-GPU prefill to 16k tokens; 80\,GB GPUs or tensor parallelism would be required for the full 128k context window. (7)~The \texttt{AdelicCache} condensation is currently validated on a dummy random-weight model; production-scale evaluation on full TinyLlama-1.1B with perplexity and NIAH retrieval benchmarks remains future work.

% ═══════════════════════════════════════════════════════════════
% 6. RELATED WORK
% ═══════════════════════════════════════════════════════════════
\section{Related Work}

\textbf{Model Surgery.} LoRA~\citep{hu2022lora} injects low-rank adapters into frozen weights for parameter-efficient fine-tuning. Our approach is complementary: rather than adapting the weight matrices, we inject a \emph{structural} modification to the attention pattern itself. The router can be composed with LoRA adapters for joint structural and parametric adaptation.

\textbf{Attention Sinks.} \citet{xiao2023efficient} identified the Attention Sink phenomenon in streaming LLMs and proposed retaining initial tokens in the KV-cache. Our work extends this finding to the block-sparse topology setting, where the sink must be anchored not just in the cache but in the \emph{routing mask} to prevent softmax collapse.

\textbf{Continuous Sparsification.} \citet{savarese2020winning} introduced continuous sparsification of neural network weights via smooth $\ell_0$ relaxations. Our Continuous Logit Homotopy adapts this principle to \emph{attention topology}: the sparsity target is not individual weight magnitudes but the block-level connectivity pattern of the attention matrix.

\textbf{Sparse Attention in Pre-Trained Models.} LongLoRA~\citep{chen2024longlora} uses shifted sparse attention with LoRA for long-context extension. MInference~\citep{jiang2024minference} identifies and exploits sparse attention patterns at inference time without retraining. Our approach differs in that the sparsity pattern is \emph{learned} during a short fine-tuning phase, not extracted post-hoc.

\textbf{Mixture of Experts.} Switch Transformer~\citep{fedus2022switch} and subsequent work~\citep{jiang2024mixtral} use learned routing for sparse feedforward computation. Our load-balancing loss follows their formulation, but applies it to \emph{attention} routing rather than expert dispatch.

\textbf{Ultrametric and $p$-adic Methods.} \citet{bradley2010mumford} connected $p$-adic analysis to hierarchical clustering. \citet{khrennikov2004padic} established foundations for $p$-adic neural networks. Our work operationalizes these mathematical structures as a differentiable inductive bias for hardware-accelerated attention.

% ═══════════════════════════════════════════════════════════════
% 7. CONCLUSION
% ═══════════════════════════════════════════════════════════════
\section{Conclusion}

Llama Surgery demonstrates that pre-trained dense language models can be continuously sparsified via differentiable topology injection, without retraining, distillation, or post-hoc pruning. The Dynamic Topology Router discovers content-based block-sparse attention patterns that are mathematically grounded in $p$-adic geometry, compatible with the Hugging Face ecosystem, and directly executable by a custom Triton kernel optimized for modern GPU architectures. When forced to perform exact sequence retrieval, the router spontaneously induces an ultrametric cophenetic hierarchy on the context window, with the multi-head architecture producing a forest ensemble of 32 independent ultrametric trees---each specializing in a distinct aspect of the retrieval task---rather than a single global hierarchy. A simulated Ring Attention deployment confirms that the emergent ultrametric topology induces consistent block-level distance separation between semantic domains, establishing the structural prerequisite for communication-aware pruning in distributed long-context inference. The \texttt{AdelicCache} condensation experiment further demonstrates that the same topological branch assignments can compress the physical KV cache from $O(N)$ to $O(W + \log N)$, maintaining a strict capacity ceiling while advancing the logical RoPE sequence position without bound---establishing the structural foundations for infinite-context generation under practical memory constraints.

The model learns to route. The kernel executes the route. The surgeon preserves the patient.

% ═══════════════════════════════════════════════════════════════
% REFERENCES
% ═══════════════════════════════════════════════════════════════
\bibliographystyle{plainnat}

\begin{thebibliography}{99}

\bibitem[Aquino-Michaels(2026)]{learningtoskip2026}
Aquino-Michaels (2026).
\newblock Learning to Skip Blocks: Self-Discovered Ultrametric Routing for Hardware-Accelerated Sparse Attention.
\newblock \emph{Preprint}.

\bibitem[Bradley(2010)]{bradley2010mumford}
Bradley, P.~E. (2010).
\newblock Mumford Dendrograms.
\newblock \emph{Computer Journal}, 53(4), 393--404.

\bibitem[Chen et~al.(2024)]{chen2024longlora}
Chen, Y., Qian, S., Tang, H., Lai, X., Liu, Z., Han, S., \& Jia, J. (2024).
\newblock LongLoRA: Efficient Fine-Tuning of Long-Context Large Language Models.
\newblock \emph{ICLR 2024}.

\bibitem[Dao et~al.(2022)]{dao2022flashattention}
Dao, T., Fu, D.~Y., Ermon, S., Rudra, A., \& R\'{e}, C. (2022).
\newblock FlashAttention: Fast and Memory-Efficient Exact Attention with IO-Awareness.
\newblock \emph{NeurIPS 2022}.

\bibitem[Fedus et~al.(2022)]{fedus2022switch}
Fedus, W., Zoph, B., \& Shazeer, N. (2022).
\newblock Switch Transformers: Scaling to Trillion Parameter Models with Simple and Efficient Sparsity.
\newblock \emph{JMLR}, 23(120), 1--39.

\bibitem[Hu et~al.(2022)]{hu2022lora}
Hu, E.~J., Shen, Y., Wallis, P., Allen-Zhu, Z., Li, Y., Wang, S., Wang, L., \& Chen, W. (2022).
\newblock LoRA: Low-Rank Adaptation of Large Language Models.
\newblock \emph{ICLR 2022}.

\bibitem[Jiang et~al.(2024a)]{jiang2024mixtral}
Jiang, A.~Q., Sablayrolles, A., Roux, A., et~al. (2024).
\newblock Mixtral of Experts.
\newblock \emph{arXiv:2401.04088}.

\bibitem[Jiang et~al.(2024b)]{jiang2024minference}
Jiang, H., Li, Y., Zhang, C., et~al. (2024).
\newblock MInference 1.0: Accelerating Pre-Filling for Long-Context LLMs via Dynamic Sparse Attention.
\newblock \emph{NeurIPS 2024}.

\bibitem[Khrennikov(2004)]{khrennikov2004padic}
Khrennikov, A.~Y. (2004).
\newblock \emph{$p$-Adic Valued Distributions in Mathematical Physics}.
\newblock Springer.

\bibitem[Merity et~al.(2017)]{merity2017pointer}
Merity, S., Xiong, C., Bradbury, J., \& Socher, R. (2017).
\newblock Pointer Sentinel Mixture Models.
\newblock \emph{ICLR 2017}.

\bibitem[Milakov \& Gimelshein(2018)]{milakov2018softmax}
Milakov, M. \& Gimelshein, N. (2018).
\newblock Online Normalizer Calculation for Softmax.
\newblock \emph{arXiv:1805.02867}.

\bibitem[Savarese et~al.(2020)]{savarese2020winning}
Savarese, P., Silva, H., \& Maire, M. (2020).
\newblock Winning the Lottery with Continuous Sparsification.
\newblock \emph{NeurIPS 2020}.

\bibitem[Xiao et~al.(2023)]{xiao2023efficient}
Xiao, G., Tian, Y., Chen, B., Han, S., \& Lewis, M. (2023).
\newblock Efficient Streaming Language Models with Attention Sinks.
\newblock \emph{ICLR 2024}.

\bibitem[Liu et~al.(2024)]{liu2024lost}
Liu, N.~F., Lin, K., Hewitt, J., Paranjape, A., Bevilacqua, M., Petroni, F., \& Liang, P. (2024).
\newblock Lost in the Middle: How Language Models Use Long Contexts.
\newblock \emph{TACL}, 12, 157--173.

\bibitem[Liu et~al.(2023)]{liu2023ring}
Liu, H., Zaharia, M., \& Abbeel, P. (2023).
\newblock Ring Attention with Blockwise Transformers for Near-Infinite Context.
\newblock \emph{arXiv:2310.01889}.

\bibitem[Zhang et~al.(2024)]{zhang2024tinyllama}
Zhang, P., Zeng, G., Wang, T., \& Lu, W. (2024).
\newblock TinyLlama: An Open-Source Small Language Model.
\newblock \emph{arXiv:2401.02385}.

\bibitem[Bolya et~al.(2022)]{bolya2022token}
Bolya, D., Fu, C.-Y., Dai, X., Zhang, P., Hoffmann, C., \& Hoffman, J. (2022).
\newblock Token Merging: Your ViT But Faster.
\newblock \emph{ICLR 2023}.

\bibitem[Bui et~al.(2026)]{bui2026make}
Bui, T., et~al. (2026).
\newblock To Make or Not to Make: The Attention Dilution Problem in Token Merging.
\newblock \emph{Preprint}.

\bibitem[Pustovit(2026)]{pustovit2026knowledge}
Pustovit, Y. (2026).
\newblock The KV Injection Dilemma: Key Arithmetic Destroys Coherence.
\newblock \emph{arXiv:2602.10822}.

\bibitem[Su et~al.(2024)]{su2024rope}
Su, J., Lu, Y., Pan, S., Murtadha, A., Bo, W., \& Liu, Y. (2024).
\newblock RoFormer: Enhanced Transformer with Rotary Position Embedding.
\newblock \emph{Neurocomputing}, 568, 127063.

\bibitem[Karpathy(2015)]{karpathy2015char}
Karpathy, A. (2015).
\newblock char-rnn.
\newblock \emph{GitHub repository}.

\end{thebibliography}

\end{document}
